\documentclass[12pt]{article}
\usepackage[T1]{fontenc}
\usepackage[utf8]{inputenc}
\usepackage{lmodern}
\usepackage{textcomp}
\usepackage{enumitem}
\usepackage{geometry}
\geometry{margin=1in}
\begin{document}
\begin{center}
Answer Key - Psychology - Undergraduate Level Assessment 22 \
\small (Answers are ordered exactly as in the question paper.)
\end{center}

\section*{Multiple Choice}
\begin{enumerate}[label=\arabic*.]
\item \textbf{Answer:} A
\\ \textit{Options:} A) "Genetic blueprint" hypothesis; B) "Innate potential" hypothesis; C) Environmental impact hypothesis; D) "Experience-driven" development theory; E) "Blank slate" hypothesis
\\ \textit{Note:} The "genetic blueprint" hypothesis posits that genetic factors play a fundamental role in shaping individual traits and behaviors, emphasizing the influence of biology in the nature-nurture controversy, which aligns with Stern's assertion that genetics significantly determine psychological characteristics.
\item \textbf{Answer:} E
\\ \textit{Options:} A) the behavior has been reinforced on a continuous schedule.; B) the behavior is not causing significant distress or dysfunction.; C) the behavior is a part of the individual's cultural or religious practices.; D) an alternative behavior cannot be identified.; E) a temporary increase in the behavior cannot be tolerated.
\\ \textit{Note:} Operant extinction may lead to an initial increase in the undesired behavior, known as an extinction burst, which could be problematic if a temporary rise in that behavior is unacceptable in the given context.
\item \textbf{Answer:} A
\\ \textit{Options:} A) Mean: 14 years, Median: 12.5, Mode: 13, Range: 8; B) Mean: 13 years, Median: 14, Mode: 12.5, Range: 17; C) Mean: 13.5 years, Median: 12.5, Mode: 14, Range: 8; D) Mean: 13 years, Median: 13, Mode: 8, Range: 9; E) Mean: 12.5 years, Median: 12, Mode: 11, Range: 10
\\ \textit{Note:} The correct answer for the measures of central tendency is derived by calculating the mean (average of the ages), median (the middle value when the ages are ordered), and mode (the most frequently occurring age), while the range is determined by subtracting the minimum age from the maximum age, highlighting the distribution and spread of the children's ages.
\item \textbf{Answer:} D
\\ \textit{Options:} A) reliability; B) face validity; C) internal validity; D) external validity; E) statistical significance
\\ \textit{Note:} External validity refers to the extent to which the findings from a sample can be generalized to a larger population, which is essential for making inferences about population behavior based on sample data.
\item \textbf{Answer:} A
\\ \textit{Options:} A) Partial-interval recording; B) Whole-interval recording; C) Duration recording; D) Time sampling recording; E) Continuous recording
\\ \textit{Note:} Partial-interval recording is a data collection method used in behavioral observation that involves counting the number of times a specific behavior occurs within predetermined time intervals, making it the correct choice for measuring occurrences of behavior over time.
\end{enumerate}

\section*{True / False}
\begin{enumerate}[label=\arabic*.]
\setcounter{enumi}{5}
\item \textbf{Answer:} False
\\ \textit{Options:} A) True; B) False
\\ \textit{Note:} A skills test may not be the most effective method for predicting job performance because it often does not account for other important factors such as personality traits, work ethic, and situational judgment, which can also significantly influence an employee's overall performance in a job.
\item \textbf{Answer:} False
\\ \textit{Options:} A) True; B) False
\\ \textit{Note:} The statement is false because sound localization primarily relies on binaural cues, such as interaural time differences and interaural level differences, rather than echo factors, loudness factors, and cognitive factors.
\item \textbf{Answer:} True
\\ \textit{Options:} A) True; B) False
\\ \textit{Note:} The statement is true because the occipital lobe is the primary visual processing center of the brain, where information received from the optic nerve is first interpreted and analyzed.
\item \textbf{Answer:} False
\\ \textit{Options:} A) True; B) False
\\ \textit{Note:} Asian and Asian-American therapy clients often express emotional problems through somatic symptoms rather than verbal symptoms due to cultural norms that emphasize physical manifestations of distress over verbal articulation of emotional issues.
\item \textbf{Answer:} True
\\ \textit{Options:} A) True; B) False
\\ \textit{Note:} The Whorfian hypothesis of linguistic relativity posits that the structure and characteristics of a language, including whether it is tonal or non-tonal, shape the way its speakers perceive and think about the world, thereby highlighting the significant differences in linguistic features across cultures.
\end{enumerate}

\section*{Long Form Response}
\begin{enumerate}[label=\arabic*.]
\setcounter{enumi}{10}
\item \textbf{Answer:} Milgram's obedience study has faced significant ethical criticisms, particularly regarding the lack of debriefing provided to participants after the experiment. This absence of debriefing meant that many participants left the study without fully understanding the nature and purpose of the experiment, potentially leading to lasting psychological distress and confusion about their actions. Ethical guidelines in psychological research emphasize the importance of informed consent and post-experimental debriefing to ensure that participants are not left with unresolved feelings of guilt or anxiety. The implications of this oversight extend beyond Milgram's study, prompting a reevaluation of research ethics in psychology, as it highlights the necessity of safeguarding participants' mental well-being and ensuring they leave the study with a clear understanding of their experience. Such considerations have since influenced the development of stricter ethical standards in psychological research.
\\ \textit{Note:} correct because it highlights the critical ethical issues surrounding Milgram's study, specifically the failure to adequately debrief participants, which undermines informed consent and may cause lasting psychological harm, thereby violating established ethical standards in psychological research.
\item \textbf{Answer:} Waiving co-payments for low-income clients within a sliding fee scale structure raises significant ethical considerations that impact both clients and insurance practices. Ethically, this approach can enhance access to necessary psychological services for those who might otherwise forego treatment due to financial barriers, promoting equity in mental health care. However, it is crucial that this waiver does not lead to profit motives from insurance companies, as this could exploit vulnerable populations. Balancing the need for affordability with the integrity of insurance practices is essential to ensure that such measures genuinely serve the interests of low-income clients without compromising the ethical standards of the psychological profession. Ultimately, while waiving co-payments can be beneficial, it must be implemented with careful consideration of its implications for the broader healthcare system.
\\ \textit{Note:} correct because it identifies the ethical balance between improving access to mental health services for low-income clients and the potential for exploitation by insurance companies, emphasizing the need for equity in healthcare while addressing financial barriers.
\end{enumerate}

\vfill
\noindent\textit{End of Answer Key}
\end{document}